% Section 1.1, from lectures/L01/appendix/a_uart_def.md, with the package itself typeset straight
% from hw/uart_def.vhd.

\section{\code{uart_def.vhd} (provided)}
\label{c1:sec:uart_def}\label{c1:app:a}
\code{uart_def} is the register map in VHDL form: one package of named constants for the seven
register indices and the meaning of each \code{STATUS}, \code{CTRL} and \code{ERROR_FLAGS} bit,
from Part~2 of the protocol (\secref{proto:part2}). Like the SPI transport, \textbf{it is given to
you}; you write every module in this course, but not the map. It exists so that the register bank
(\code{uart_regs}, \chapref{5}) and the testbenches can name registers and bits instead of writing
bare numbers, and so that the two sides of the wire cannot drift apart.

The C++ driver carries the same map in \file{register_map.hpp}, transcribed once on each side of the
wire. Both store \textbf{bit positions}, not masks, so \code{STATUS} bit~1 is the constant \code{1}
in both languages. The C++ side declares only the subset the driver actually uses, under the
\code{reg::}, \code{status::}, \code{ctrl::} and \code{error::} namespaces; where a name exists on
both sides, the value is identical. You do not edit \code{uart_def}, but read it closely: a module
that indexes the wrong bit is a bug you would only find on the bench.

\subsection{The constants}
\label{c1:sec:constants}
Everything is a \code{natural} in \code{package uart_def}. The register indices are the offset
divided by 4; the bit constants are positions into a 32-bit register word.

\begin{booktable}{@{}p{7.4em}p{8.6em}rL@{}}
  \thead{Group} & \thead{Constant} & \thead{Value} & \thead{Meaning} \\
  \midrule
  Register index & \code{REG_STATUS} & 0 & Status register. \\
   & \code{REG_CTRL} & 1 & Control register. \\
   & \code{REG_BAUD_DIV} & 2 & Baud divider. \\
   & \code{REG_TX_DATA} & 3 & Push a byte into the TX FIFO. \\
   & \code{REG_RX_DATA} & 4 & Front byte of the RX FIFO (a pure read). \\
   & \code{REG_RX_POP} & 5 & Advance the RX FIFO. \\
   & \code{REG_ERR_FLAGS} & 6 & Error flag register. \\
  \midrule
  \code{STATUS} bit & \code{ST_TX_READY} & 0 & The TX FIFO is not full. \\
   & \code{ST_RX_VALID} & 1 & The RX FIFO is not empty. \\
   & \code{ST_ERROR} & 2 & One or more error flags are set. \\
   & \code{ST_TX_IDLE} & 3 & The TX FIFO is empty and the line is idle. \\
  \midrule
  \code{CTRL} bit & \code{CT_ENABLE} & 0 & Enable the peripheral. \\
   & \code{CT_PARITY_LO} & 1 & Parity select, low bit. \\
   & \code{CT_PARITY_HI} & 2 & Parity select, high bit (00 none, 01 even, 10 odd). \\
   & \code{CT_STOP} & 3 & 0 for 1 stop bit, 1 for 2 stop bits. \\
   & \code{CT_RX_IRQ} & 4 & RX-valid IRQ mask. \\
   & \code{CT_TX_IRQ} & 5 & TX-ready IRQ mask. \\
  \midrule
  \code{ERROR_FLAGS} bit & \code{ER_FRAMING} & 0 & Framing error. \\
   & \code{ER_PARITY} & 1 & Parity error. \\
   & \code{ER_OVERRUN} & 2 & Overrun. \\
\end{booktable}

Because they are positions, you index a \code{std_logic_vector} with them directly, for example
\code{status(ST_RX_VALID)} for the RX-valid bit, exactly as \code{1U << status::RX_VALID} forms the
same mask on the C++ side. Laid out as words, the three flag registers are shown in
Figure~\ref{c1:fig:reg_fields}.

\bookfigure{reg_fields}{The bit fields of \code{STATUS}, \code{CTRL} and \code{ERROR_FLAGS}}
  {c1:fig:reg_fields}

\code{CTRL} is the only one with bits this course does not act on, and \code{ERROR_FLAGS} the only
one where two of the three have no producer yet; both are noted where they are built, in
\secref{c5:app:a}.

\subsection{The \code{to_hex} helper}
\label{c1:sec:to_hex}
The package also carries one small convenience for the testbenches: \code{to_hex}, which renders a
\code{std_logic_vector} as an uppercase, zero-padded hex string. It pads to whole nibbles taken from
the argument's own length, so \code{x"A5"} becomes \code{"A5"} and a 12-bit bus becomes three
digits. It lets a testbench name the offending value in a failure message,

\begin{vhdlcode}
report "rx_data byte 1 wrong: expected 0xA5, got 0x" & to_hex(captured_rx) & "!"
    severity error;
\end{vhdlcode}

\noindent rather than a bare ``mismatch''. It is the one piece of behaviour in an otherwise pure
constants package: declared in the package, defined in the package body, and called only by
testbenches. No synthesizable module uses it.

The whole package, constants and helper, is short enough to print as it stands in the repository:

\vhdlfile{hw/uart_def.vhd}

\subsection{Where it fits}
\label{c1:sec:uart_def_fits}
\code{uart_def} sits underneath everything that speaks in registers. \code{uart_regs}
(\chapref{5}) is the module that actually reads and writes these bits; the package only names them.
\code{uart_top} never uses it at all, because the top wires the register bus without interpreting
anything carried on it. Six testbenches analyze against the package, most of them only for
\code{to_hex}.

On the far side of the SPI wire, the driver you build from \chapref{6} on reaches the same
registers through the same positions. That is the point of transcribing the map twice rather than
inventing it twice: the protocol is the single source, and each language is checked against it
rather than against the other.

\subsection{What's ahead}
\secref{c1:sec:uart_top} is the peripheral top, \code{uart_top}, the structural skeleton built
against this package. The exercises in \secref{c1:sec:exercises} read \code{uart_def}, then build
\code{uart_top} from the ground up. The datapath and register modules the top composes arrive in
Chapters~2 to~5, and are instantiated into \code{uart_top} as each one lands.
